\scanpage{118}
\mspage{105}
\[
 X_U=\lim_{i\in U}X(i),\qquad X_V=\underset{i\in V}{\unsure{\underrightarrow{\lim}|\underline{\lim}}}\,X(i),
\]
et de $=$ pour $Y$, et considérons le
diagramme

\diagram{\begin{tikzcd}[ampersand replacement=\&,column sep=large]
\& X_U \arrow[r,"\text{trans.}"] \arrow[d,"f_U"'] \arrow[dr,phantom,"\text{cocart}"] \& X_V \arrow[d] \arrow[dd,bend left=48,dashed,"f_V"]\\
(6.2.1) \& Y_U \arrow[r,"{\deleted{\unsure{\text{trans.}}}}"'] \arrow[dr,dashed,"\text{trans.}"'] \& Y_{U,V} \arrow[d,"{f_{U,V}}"']\\
\& \& Y_V
\end{tikzcd}}

Alors la flèche $u_{U,V}$ est dans $C$ (resp. $TC$).

Démonstration de 6.2. (\uuline{Admettons} \deleted{$U,V$}
que la restriction d'une flèche de
$\deleted{\unsure{f}} C(I)$ resp $TC(I)$ à un fermé
\deleted{\unsure{de}} $\deleted{\unsure{U}}V$ de $I$ \deleted{de} est dans $\deleted{\unsure{f}} C(V)$ \unsure{respt.}
$TC(V)$ Cf. \ldots\ ci-dessous). Cela
nous permet de supposer $V=I$.

La donnée de $U$, qui est un ouvert
de $\unsure{I|I^\circ}$, équivaut à la donnée d'une
flèche \unsure{$v$|``} $I^\circ\longrightarrow\Delta^1$ (\deleted{\unsure{par}} prenant la
fibre en $0$), ou encore à la donnée
d'une flèche $u\colon I\unsure{\longrightarrow|\text{---}}(\Delta^1)^\circ\simeq\Delta^1$
(en prenant la fibre en $1$), \unsure{soit}
\[
 u\colon I\longrightarrow\Delta^1,\qquad
 U=\deleted{\unsure{V}}I_1\overset{\mathrm{def}}{=}u^{-1}(\{1\}).\tag*{(6.2.2)}
\]
% archive sheet number: 117

